\section{Requirements and Threat Model}

\subsection{Functional}
Sub-kB/day, intermittent links, device $\rightarrow$ gateway verification, daily public timestamp proof.
Anti-replay windowed acceptance at the gateway (default $w{=}64$), deterministic/canonical hashing for reproducibility.

\subsection{Non-functional}
Ultra--low-power MCU, simple provisioning, reproducible hashing.

% --- Insert: Project scope and roadmap ---
\subsection{Project scope and roadmap}
TrackOne is explicitly framed as ``Track One of Three'': the core pipeline that provides ingestion, cryptographic integrity, deterministic batching, and public anchoring. It is an intentionally narrow, verifiability-first track whose deliverables are:
\begin{itemize}[nosep]
  \item A hardened ingestion pipeline (framed NDJSON + AEAD semantics) and replay protection;
  \item Deterministic canonicalization and Merkle batching with day-chunked artifacts;
  \item Public anchoring (OpenTimestamps) and verifier tooling that lets third parties recompute and attest daily roots;
  \item Operational guidance and minimal hardware reference that satisfy the digital contract (frame shape, nonce construction, provisioning).
\end{itemize}

Features reserved for later tracks (Track Two / Three) include field-deployable control loops, sophisticated on-pod actuation safety interlocks, large-scale analytics and UI, and long-running automated key-rotation/PKI operations. Position TrackOne as an MVP that is ``production-informed'' but not yet a full safety-certified control system.

% --- Insert: Operational and environmental constraints ---
\subsection{Operational and environmental constraints}
For subterranean irrigation deployments, the physical environment and service model materially change requirements. We therefore set explicit operational constraints that must be observed when deploying TrackOne artifacts in the field:
\begin{itemize}[nosep]
  \item Deployment tolerance: pods are expected to be wet, buried, and inaccessible for long durations; hardware must therefore favour robustness and conservative scheduling for maintenance windows.
  \item Maximum allowed downtime (gateway+backend) for a site: configurable, typical value 24 hours; longer outages must be escalated to human operators.
  \item Acceptable data loss: define a site-specific bound (e.g., no more than 0.5\% of day's facts lost) beyond which operator rollback or manual audit is required.
  \item Sampling frequency: default 1--10 samples per hour per pod depending on sensor; higher rates must be justified by an explicit risk/benefit analysis.
  \item Firmware update policy: updates must support atomic rollback and require an operator-confirmed maintenance window; over-the-air auto-update that can change actuation logic is out-of-scope for TrackOne and deferred to Track Two.
  \item Storage limits: on-device and gateway storage quotas must be set; when storage is full, pods/gateway must stop non-essential writes and raise an alert to operator channels rather than silently dropping telemetry.
\end{itemize}

% --- strengthen threat model ---
\subsection{Threat model}
Passive RF eavesdropping; active injection/replay; rogue gateway; physical capture (NVM readout). Additionally, for irrigation deployments we explicitly consider:
\begin{itemize}[nosep]
  \item Communication loss: prolonged uplink outages and partial replication windows; consequences include delayed anchoring and postponed verification.
  \item Clock drift / time sync failures: incorrect timestamps can corrupt day-chunk boundaries and complicate forensic recomputation.
  \item Loss/corruption of field logs: storage device failures or accidental truncation of \texttt{out/} artifacts.
  \item Unsafe actuation: if the system is later used to drive valves/relays, corrupted inputs or stale data may cause over- or under-irrigation with real-world consequences.
\end{itemize}

Out of scope: side-channels; jamming. Mitigations: framed anti-replay (ADR-002); deterministic artifacts and public timestamping (ADR-003); \textbf{peer attestation of day roots} (Ed25519 quorum) to reduce single-gateway trust; optional \textbf{parallel RFC~3161 TSA anchoring} (ADR-015) alongside OTS for PKI-backed timing.

% --- Insert: irrigation-specific requirements (concise, testable) ---
\subsection{Irrigation-specific requirements (traceable)}
To support traceability and operational safety in irrigation contexts, we add a small set of testable, measurable requirements that must be in scope for any deployment using TrackOne artifacts:
\begin{itemize}[nosep]
  \item R-1 (Maximum acceptability of delayed anchoring): For any given day, anchors (OTS proofs) must be present within 72 hours of day end or the system must mark the day as ``pending anchor'' and notify operators.
  \item R-2 (Maximum tolerable data loss): Daily fact loss must not exceed X\% (site-defined; default 0.5\%). Exceeding this triggers operator review and a preservation workflow.
  \item R-3 (Sampling \& action separation): Control decisions that affect irrigation volumes MUST not be taken on facts that are older than T hours (default: 24 h) without human confirmation.
  \item R-4 (Update rollback): Any firmware update that modifies actuation logic must support a verified rollback mechanism and be staged to a small pilot group for at least 7 days before full rollout.
  \item R-5 (Auditability): All critical operator actions (e.g., manual valve overrides, firmware promote) must be logged with operator identity, timestamp, and Merkle-anchored proof of the day's root that covers the action metadata.
\end{itemize}
